% Phase III status addendum — appended after phase2_findings_brief.pdf.
% Documents the Route D negative finding and the consequences for the
% open closed-form problem.

\documentclass[11pt]{article}
\usepackage[margin=0.92in]{geometry}
\usepackage[T1]{fontenc}
\usepackage[utf8]{inputenc}
\usepackage{lmodern}
\usepackage{microtype}
\usepackage{amsmath,amssymb,amsthm,mathtools,bm}
\usepackage{booktabs,array,tabularx,longtable}
\usepackage{enumitem}
\usepackage[dvipsnames]{xcolor}
\usepackage{hyperref}
\usepackage{fancyhdr}
\usepackage{titlesec}

\definecolor{TitleBlue}{HTML}{173753}
\definecolor{AccentBlue}{HTML}{2F5D7E}
\definecolor{SoftBlue}{HTML}{EEF4FA}
\definecolor{SoftGold}{HTML}{FAF6EE}
\definecolor{SoftGreen}{HTML}{EEF8F0}
\definecolor{SoftRed}{HTML}{FCEEEE}
\definecolor{RuleGray}{HTML}{D7E1EA}
\definecolor{MutedText}{HTML}{4E5A66}

\hypersetup{colorlinks=true,linkcolor=AccentBlue,urlcolor=AccentBlue,
  citecolor=AccentBlue,
  pdftitle={Phase III status — intrinsic U(3) richness of the inter-window matrix},
  pdfauthor={Type-1 N=3 Landau-Zener project}}
\pagestyle{fancy}\fancyhf{}
\fancyhead[L]{\textcolor{MutedText}{\small Phase III status}}
\fancyhead[R]{\textcolor{MutedText}{\small Type-1, $N=3$ Landau--Zener}}
\fancyfoot[C]{\textcolor{MutedText}{\thepage}}
\setlength{\headheight}{14pt}\renewcommand{\headrulewidth}{0.4pt}
\titleformat{\section}{\large\bfseries\color{TitleBlue}}{\thesection}{0.6em}{}
\titleformat{\subsection}{\normalsize\bfseries\color{AccentBlue}}{\thesubsection}{0.5em}{}
\setlength{\parindent}{0pt}\setlength{\parskip}{0.5em}
\setlist[itemize]{leftmargin=1.4em,itemsep=0.18em,topsep=0.18em}

\newcommand{\C}{\mathbb C}\newcommand{\R}{\mathbb R}
\newcommand{\ii}{\mathrm i}\newcommand{\dd}{\mathrm d}
\newcommand{\im}{\operatorname{Im}}
\newcommand{\re}{\operatorname{Re}}

\newcommand{\infobox}[2]{\begin{center}\fcolorbox{RuleGray}{#1}{\begin{minipage}{0.94\linewidth}\small #2\end{minipage}}\end{center}}

\begin{document}

\begin{center}
{\Large\color{AccentBlue}\bfseries Phase III status\par}\vspace{0.45cm}
{\fontsize{19}{23}\selectfont\bfseries\color{TitleBlue}Intrinsic $\mathrm U(3)$ richness of the\par}
{\fontsize{19}{23}\selectfont\bfseries\color{TitleBlue}inter-window matrix\par}
\vspace{0.30cm}
{\large\color{MutedText}Route D negative result and consequences for Routes A--E\par}
\end{center}

\infobox{SoftBlue}{%
\textbf{Phase III headline.} The Phase II finding that the inter-window
matrix $D \in \mathrm U(3)$ carries irreducibly rich structure has been
sharpened: \emph{the structure persists in any LZ-block convention.}
Re-extracting $D$ using the proper Stokes-phased Landau--Zener amplitude
blocks $N_X^{\text{proper}}$ (with $\varphi_S(\delta) = \delta(\log
\delta - 1) + \pi/4 + \arg\Gamma(1-\ii\delta)$) does \emph{not} reduce
the residual of the $(\mathrm{diag}\cdot R_{(\mathrm{sp}_A, \mathrm{sp}_B)})$
ansatz: median residual on a 12-sample calibration is $0.0664$ in both
the phase-free and the Stokes-phased convention, identical to within
six significant figures, across all four sign-convention combinations
$(\sigma_A, \sigma_B) \in \{\pm 1\}^2$. The Phase~III plan's
falsifiable exit ``Route D fails: the LZ block convention is not the
issue; the $\mathrm U(3)$ richness is intrinsic'' is therefore
\emph{triggered}.}

\section{Route D execution and outcome}
\label{sec:route-d}

Route D was the foundation step of the Phase III plan: re-extract the
empirical inter-window matrix
$D = N_B^{-1} U_{\text{bench}} N_A^{-1}$ using Stokes-phased LZ
amplitude blocks rather than the phase-free real-orthogonal block of
\texttt{assay.closed\_form.lz\_rotation}. The proper block carries
$\varphi_S(\delta)$ on its diagonal entries:
\[
N_{\text{proper}}(p,\delta) \;=\;
\begin{pmatrix}
\sqrt{1-p}\,\mathrm e^{-\ii \sigma\,\varphi_S(\delta)} & \sqrt p \\
-\sqrt p & \sqrt{1-p}\,\mathrm e^{+\ii \sigma\,\varphi_S(\delta)}
\end{pmatrix},
\]
with $\sigma \in \{+1, -1\}$ the convention's sign degree of freedom.

\subsection{Calibration result}

On a 12-sample calibration subset (3 each from the four success strata
$\{S1, S6, S7, S8\}$), all four $(\sigma_A, \sigma_B)$ combinations
produce the \emph{identical} median residual of the
$(\mathrm{diag}\cdot R)$ ansatz reconstruction:

\begin{center}\small
\begin{tabular}{lccc}
\toprule
$(\sigma_A, \sigma_B)$ & median resid & max resid & wall time \\
\midrule
$(+1, +1)$ & $0.0664$ & $0.7681$ & $90\,$s \\
$(+1, -1)$ & $0.0664$ & $0.7681$ & $84\,$s \\
$(-1, +1)$ & $0.0665$ & $0.7681$ & $102\,$s \\
$(-1, -1)$ & $0.0665$ & $0.7681$ & $91\,$s \\
\midrule
phase-free baseline & $0.0664$ & $1.4454$ & $\sim 90\,$s \\
\bottomrule
\end{tabular}
\end{center}

The Stokes-phased convention modestly reduces the \emph{tail}
(max-residual $1.4454 \to 0.7681$, a $2\times$ improvement at the
worst samples) but does not move the \emph{median} at all. Whatever
phase content the Stokes-phased $N_X^{\text{proper}}$ captures, it is
not the dominant source of $D$'s rank-three $\mathrm U(3)$ richness.

\subsection{The intrinsic off-block content}

Direct inspection of a worst-case S2 sample with intermediate
$(p_A, p_B) \approx (0.43, 0.02)$ shows that the residual of the
$(\mathrm{diag}\cdot R_{(\mathrm{sp}_A,\mathrm{sp}_B)})$ ansatz lies
\emph{outside} the spectator $(\mathrm{sp}_A, \mathrm{sp}_B)$ plane:

\[
D \;\approx\; \begin{pmatrix}
+0.51+0.19\ii & +0.32+0.15\ii & +0.12+0.75\ii \\
+0.02-0.56\ii & -0.10+0.82\ii & -0.02+0.01\ii \\
-0.36+0.51\ii & -0.28+0.33\ii & +0.65+0.07\ii
\end{pmatrix}
\quad
(D - D_{\mathrm{model}})_{\mathrm{im}} \;\approx\; \begin{pmatrix}
-0.16 & +0.15 & +0.75 \\
-0.56 & -0.12 & -0.28 \\
+0.51 & +0.36 & -0.03
\end{pmatrix}.
\]

The model residual entries at positions $(0,1)$, $(0,2)$, $(1,0)$,
$(2,0)$, $(1,2)$, $(2,1)$ are all of order $0.3$--$0.8$. The
$(\mathrm{diag}\cdot R)$ ansatz forces the common active index
$i_c = 0$ to be a fixed point of the rotation, so its row and column
should be diagonal + zeros. They are not: $D[0,1]$, $D[0,2]$,
$D[1,0]$, $D[2,0]$ all carry $O(1)$ magnitudes.

\textbf{The closed form must produce a full $\mathrm U(3)$, not a
$\mathrm U(2)\oplus\mathrm U(1)$ subgroup.}

\section{Consequences for Routes A, B, C}

The Phase~III plan's risk-ranked outlook had Route C (AKT/DDP) as the
highest-confidence \emph{predictive} route on the grounds that its
architecture had the right spectator-plane rotation. The empirical
$D$ now invalidates this architecture: a rotation in the
$(\mathrm{sp}_A, \mathrm{sp}_B)$ plane is structurally a
$\mathrm U(2)\oplus\mathrm U(1)$ element, but $D$ is a full
$\mathrm U(3)$ element. \emph{Route C as architecturally
specified cannot reach the closed form, regardless of how well its
Bessel block is implemented.}

The implementation work executed in Phase~III confirms the failure
modes empirically:

\begin{itemize}
 \item \textbf{Route C extension (\texttt{formula\_DDP\_aki}).} Replaced
 the $(\mathrm{sp}_A, \mathrm{sp}_B)$ SO(2) rotation with a modified-Bessel
 block at index $\nu = 1/3$, with the corrected
 $V_{AB} = \int(L_H/p)\sqrt{Q_4}\dd u$ Voros symbol on the spinor cover.
 Smoke test on 40 samples: per-stratum mean residual $0.60$--$0.82$,
 \emph{worse} than the original \texttt{formula\_DDP}'s $0.43$--$0.60$.
 Two causes: (i) the implemented Bessel block uses a heuristic
 parameterisation ($\sin\theta = \sin(\arg V)$ for the mixing magnitude,
 not the proper Wronskian-normalised $K_\nu/I_\nu$ ratio); (ii) the
 ``corrected'' $V_{AB}$ on the exterior sheet $\lambda_0$ is exactly
 $2\times$ the old $V$ across all samples, suggesting the integration
 contour is on the wrong sheet.

 \item \textbf{Route A extension
 (\texttt{formula\_HN\_nonabelian}).} Wired the existing
 \texttt{spectral\_network.py} machinery into the connection-matrix
 product via wall-jump matrices $K_\gamma = I + \mu_\gamma E_{ij}$ at
 real-axis wall crossings. The implementation overflows
 catastrophically (per-stratum residual $10^{23}$ to $10^{193}$).
 Cause: the Voros symbol $\mu = \mathrm e^{\ii I_X}$ with
 $\im I_X < 0$ produces $|\mu| = \mathrm e^{|\im I_X|}$ which is huge.
 The naive unipotent factor $I + \mu E_{ij}$ is non-unitary; without
 careful side-tracking (pairing each wall jump with a conjugate factor
 whose product is unitary), the product loses unitarity by hundreds
 of orders of magnitude.

 \item \textbf{Route B extension.} Not implemented in Phase~III, but
 the log-space + Faddeev finite-$b$ refactor remains the prescribed
 path and is a candidate for future work. Its closed-form image is
 a full $\mathrm U(3)$ via the mutation product, so it is not
 architecturally excluded by Route D's finding.
\end{itemize}

\section{Refined open problem}

The closed-form challenge is now precisely:

\infobox{SoftGold}{%
\textbf{Find a $\mathrm U(3)$-valued function
$D(\epsilon, \gamma, a)$} whose $|\cdot|^2$ projection plus row/column
stochasticity gives the two remaining transcendental entries of $P$.
The function must:

\begin{itemize}
 \item be parametrised by 8 real parameters (6 after the global U(1)
 gauge), \emph{not} reducible to a 2-plane rotation;
 \item match the empirical $D$ on 300 stratified samples to U(3)
 geodesic distance $\le 10^{-8}$;
 \item be derivable from one of the three exact-WKB literature
 frameworks (Hollands--Neitzke, Iwaki--Nakanishi, AKT) OR surfaced
 directly from data via symbolic regression.
\end{itemize}}

This is a research-level problem. The Phase~II/III prototypes
collectively demonstrate that none of the standard sequential 2-level LZ
products, restricted spectator-rotation ansätze, or single-window
Stokes-phase corrections can produce it.

\section{Recommended next steps}

\begin{enumerate}
 \item \textbf{Proper Hollands--Neitzke side-tracking (Route A
 done correctly).} The \texttt{spectral\_network} module identifies
 all real-axis wall crossings, BPS saddles, and sheet pairs. The
 missing piece is the GMN ``$\omega$-pair'' construction: each BPS
 state contributes a pair of unipotent factors with reciprocal Voros
 symbols whose product is a unitary $2\times 2$ rotation. Implementing
 this side-tracking is mechanical but requires careful attention to
 which side of each wall each chamber sits on. Estimated effort:
 5--10 days.
 \item \textbf{Iwaki--Nakanishi with log-space arithmetic and
 finite-$b$ Faddeev (Route B done correctly).} Produces a
 full $\mathrm U(3)$ via the cluster mutation product. The
 implementation pieces (log-space matrix accumulation; finite-$b$
 Faddeev convergence) are well-defined. Estimated effort: 3--5 days.
 \item \textbf{PySR symbolic regression on the empirical $D$
 (Route E with proper backend).} Install PySR (requires the Julia
 toolchain, an explicit permission grant needed in this auto-mode
 environment), feed it the 300-sample $D$ dataset, and let it search
 over an operator grammar that includes $\Gamma$-functions, modified
 Bessel functions, and quantum dilogarithms. If the closed form is a
 finite product of named transcendentals at parameter-determined
 arguments (the expected shape), PySR will surface it.
 \item \textbf{Deep literature dive on Hollands--Neitzke
 non-abelianisation at the genus-1 Hitchin curve.} The Type-1, $N=3$
 spectral curve has $b_1(\Sigma_Y) = 2$, matching the window count.
 The Hollands--Neitzke 2019 paper on non-abelianisation for
 genus-1 Hitchin systems is the direct reference. A careful read
 should resolve the side-tracking and $\theta$-phase conventions
 that defeated the Phase~III prototype.
\end{enumerate}

\section{Status of the closed-form push}

\begin{itemize}
 \item \textbf{Exact \& elementary} (Phase~I deliverables): two
 Brundobler--Elser extreme-survival entries. Validated to numerical
 floor on 300 samples.
 \item \textbf{Empirically extracted} (Phase~III deliverable): the
 inter-window matrix $D \in \mathrm U(3)$ on 300 samples, in two
 conventions (phase-free, Stokes-phased), with the
 $(\mathrm{diag}\cdot R)$ ansatz residual quantified.
 \item \textbf{Constrained} (Phase~III finding): $D$ is structurally
 a full $\mathrm U(3)$, not a $\mathrm U(2)\oplus\mathrm U(1)$
 subgroup. Architectural ansätze that produce only a spectator-plane
 rotation are excluded.
 \item \textbf{Open}: the analytic form of
 $D(\epsilon, \gamma, a)\in\mathrm U(3)$.
\end{itemize}

\section{Artifacts}

\textbf{New code (Phase~III):}
\begin{itemize}
 \item \texttt{assay/lz\_stokes\_phased.py} -- Stokes-phased LZ amplitude block.
 \item \texttt{assay/d\_extract.py} -- shared $D$-extraction primitives,
 U(3) geodesic distance.
 \item \texttt{assay/convergence\_test.py} -- pairwise route comparison
 and empirical $D$ agreement gate.
 \item \texttt{notes/phase2\_D\_stokes\_phased.py} -- Route D driver
 (sign calibration + 300-sample sweep).
 \item \texttt{notes/phase2\_symbolic\_regression\_lite.py} -- numpy-only
 fallback when PySR is unavailable.
 \item \texttt{assay/ddp.py} (extended) -- \texttt{formula\_DDP\_aki},
 \texttt{aki\_bessel\_block}, \texttt{inter\_window\_action\_corrected},
 \texttt{ddp\_D}.
 \item \texttt{assay/hitchin\_holonomy.py} (extended) --
 \texttt{formula\_HN\_nonabelian}, \texttt{hn\_connection\_matrix\_nonabelian},
 \texttt{hn\_D}, \texttt{\_wall\_data}.
\end{itemize}

\textbf{Output datasets:}
\begin{itemize}
 \item \texttt{notes/phase2\_D\_stokes\_phased.csv} -- 300 samples × 9
 complex $D$ entries × (theta, chi, model\_resid, ...) in Stokes-phased
 convention.
 \item \texttt{assay/empirical\_D\_cache.pkl} -- per-record cache for
 \texttt{empirical\_D}.
\end{itemize}

\textbf{Plan file:} \texttt{vast-drifting-avalanche.md} (the approved
Phase~III plan).

\end{document}
